\chapter{Brain Aneurysm}
\index{Brain Aneurysm}

\section*{Definition and Overview}
A brain (cerebral or intracranial) aneurysm is an abnormal, balloon-like bulge that develops in the wall of an artery inside the skull, most often at a point where the vessel branches. The arterial wall weakens in a focal area and, under the continuous stress of blood flow, gradually distends outward into a thin-walled pouch. The great majority of brain aneurysms never rupture, produce no symptoms, and are discovered incidentally on imaging performed for another reason. The principal danger is rupture: a burst aneurysm bleeds into the space between the brain and its surrounding membranes -- a subarachnoid haemorrhage -- which is a sudden, life-threatening event requiring immediate emergency care. Some aneurysms are managed by monitoring and risk-factor control alone, while others are treated proactively with endovascular or surgical techniques to prevent rupture.

\section*{Epidemiology}
Unruptured intracranial aneurysms are relatively common, found in roughly 3--5\% of the adult population, although most never rupture. They are slightly more common in women than in men, and prevalence rises with age, with most aneurysms detected between the ages of 40 and 60. Ruptured aneurysms are the leading cause of non-traumatic subarachnoid haemorrhage, accounting for about 85\% of such bleeds; the overall incidence of aneurysmal subarachnoid haemorrhage is around 6--10 cases per 100,000 people per year. Certain groups are at markedly higher risk, including people with two or more affected first-degree relatives, those with autosomal dominant polycystic kidney disease, and those with connective tissue disorders such as Marfan syndrome, Ehlers--Danlos syndrome, and neurofibromatosis type 1. The incidence of rupture has been declining in recent decades in developed countries, in part because of better blood-pressure control and lower smoking rates.

\section*{Aetiology and Pathophysiology}
Aneurysms form where the muscular wall of an artery is congenitally or acquiredly weakened, typically at bifurcations, where the shear stress of turbulent blood flow is greatest. The wall bulges, and the aneurysm enlarges slowly over years. Several factors promote formation and rupture:

\begin{itemize}
    \item \textbf{Hypertension:} the most important modifiable risk factor, placing continuous strain on the vessel wall
    \item \textbf{Smoking:} damages the endothelium and accelerates both aneurysm formation and growth
    \item \textbf{Excessive alcohol consumption:} raises blood pressure and is associated with a higher rupture risk
    \item \textbf{Family history:} a first-degree relative with an aneurysm approximately doubles the risk
    \item \textbf{Inherited conditions:} including autosomal dominant polycystic kidney disease and connective tissue disorders
    \item \textbf{Stimulant drugs:} particularly cocaine, which causes sudden hypertensive surges
    \item \textbf{Infection or trauma:} rarely produces so-called mycotic or traumatic aneurysms
\end{itemize}

When an aneurysm ruptures, arterial blood escapes into the subarachnoid space under high pressure. The sudden accumulation of blood raises intracranial pressure dramatically, which may transiently reduce cerebral blood flow, and the irritant effect of blood provokes intense meningeal inflammation and can trigger vasospasm -- a dangerous narrowing of cerebral arteries that typically peaks four to ten days later and can cause delayed ischaemic stroke. Rebleeding, hydrocephalus, and disturbances of salt and water balance are further consequences of the haemorrhage.

\section*{Clinical Features}
An unruptured aneurysm is usually silent. When symptoms do occur, they are typically caused by pressure on adjacent structures, especially in larger lesions:

\begin{itemize}
    \item A drooping eyelid or double vision, from compression of the oculomotor nerve
    \item Pain behind or above the eye
    \item A dilated pupil or facial numbness
    \item Occasionally a persistent headache
\end{itemize}

The hallmark of rupture is the sudden onset of an explosive headache, classically described as the worst headache of the patient's life and reaching maximal intensity within seconds to minutes (a thunderclap headache). It is often accompanied by:

\begin{itemize}
    \item Neck stiffness and sensitivity to light (photophobia)
    \item Nausea and vomiting
    \item Transient or prolonged loss of consciousness
    \item Confusion or agitation
    \item Seizures
    \item Focal neurological deficits such as weakness or difficulty speaking
\end{itemize}

Not all thunderclap headaches are caused by aneurysms, but any sudden severe headache demands urgent evaluation to exclude one.

\section*{Differential Diagnosis}
The presentation of a ruptured aneurysm overlaps with several other conditions:

\begin{itemize}
    \item Thunderclap headache from benign causes, such as cough headache or orgasmic headache
    \item Migraine, including hemiplegic migraine, which can mimic focal deficits
    \item Meningitis -- inflammation of the meninges producing similar meningism
    \item Ischaemic or haemorrhagic stroke
    \item Cervical artery dissection, which can cause sudden headache and neurological signs
    \item Reversible cerebral vasoconstriction syndrome
    \item Pituitary apoplexy or bleeding into a brain tumour
\end{itemize}

\section*{When to Seek Medical Care}
A sudden, severe, explosive headache is a medical emergency even if it settles, and it should never be dismissed or treated at home. If you have been diagnosed with an aneurysm and develop a new headache, double vision, or any neurological symptom, seek assessment without delay.

\textbf{Contact your local emergency services for:}
\begin{itemize}
    \item A sudden, explosive headache -- the worst headache of your life
    \item Collapse or loss of consciousness
    \item Any seizure or fit
    \item Sudden weakness, numbness, difficulty speaking, or confusion
    \item Neck stiffness together with a severe headache and fever
    \item A drooping eyelid, double vision, or a dilated pupil that appears suddenly
\end{itemize}

\section*{Diagnosis and Investigations}
\begin{itemize}
    \item \textbf{CT scan of the head:} the first-line investigation for suspected subarachnoid haemorrhage; it is highly sensitive within the first six to twelve hours and becomes less sensitive over subsequent days
    \item \textbf{Lumbar puncture:} if the CT is normal but the clinical picture strongly suggests a bleed, examination of the cerebrospinal fluid for blood and xanthochromia may confirm the diagnosis
    \item \textbf{CT angiography or MR angiography:} non-invasive imaging of the cerebral vessels to identify and characterise the aneurysm, its size, and its position
    \item \textbf{Catheter (digital subtraction) angiography:} the most detailed vascular imaging, used to plan endovascular or surgical treatment and to detect additional small aneurysms
    \item \textbf{Screening imaging:} offered to people with a strong family history or predisposing inherited conditions, so that an unruptured aneurysm can be detected before it ruptures
    \item \textbf{Baseline and serial assessment:} for conservatively managed lesions, regular follow-up imaging to monitor for growth
\end{itemize}

\section*{Management}
Treatment is individualised and depends on whether the aneurysm has ruptured, its size, site, morphology, and the patient's age and general health:

\begin{itemize}
    \item \textbf{Unruptured aneurysm:} small lesions with a low rupture risk may be monitored with serial scans and rigorous risk-factor control; larger or higher-risk lesions are generally treated to prevent rupture, weighing the natural history against the risks of intervention
    \item \textbf{Endovascular coiling:} platinum coils (and, increasingly, flow-diverting stents) are deployed through a catheter advanced from the groin to occlude the aneurysm from the circulation
    \item \textbf{Surgical clipping:} a craniotomy is performed and a small metal clip is placed across the neck of the aneurysm to exclude it from blood flow
    \item \textbf{After rupture:} emergency stabilisation, careful blood-pressure management, and early treatment to secure the aneurysm, since the risk of fatal rebleeding is highest in the first days
    \item \textbf{Medication:} nimodipine is given after rupture to reduce the risk of delayed ischaemia from vasospasm; analgesia and antiemetics are used for symptom control
    \item \textbf{Complication management:} treatment of hydrocephalus by external drainage, correction of hyponatraemia, and management of seizures
    \item \textbf{Rehabilitation:} physiotherapy, speech and occupational therapy, and neuropsychological support for survivors of rupture
\end{itemize}

\section*{Complications}
\begin{itemize}
    \item Rebleeding, which carries a very high risk of death or severe disability
    \item Cerebral vasospasm and delayed ischaemic stroke in the days to weeks after rupture
    \item Hydrocephalus from blood obstructing cerebrospinal fluid flow
    \item Hyponatraemia due to cerebral salt wasting or the syndrome of inappropriate antidiuretic hormone secretion
    \item Seizures and late epilepsy
    \item Cognitive, memory, and mood disturbance, which may be substantial after a bleed
    \item Complications of treatment itself, including procedural stroke, aneurysm rupture, or recurrence requiring further treatment
\end{itemize}

\section*{Prognosis and Outcomes}
The outlook depends critically on whether the aneurysm has ruptured. Unruptured aneurysms carry a low annual risk of rupture -- generally well under 1\% per year for small lesions in the anterior circulation -- and many are safely monitored for years. Once rupture occurs, however, the condition is grave: approximately 40--50\% of people who experience an aneurysmal subarachnoid haemorrhage die, and many survivors have permanent neurological or cognitive disability. Outcomes have nonetheless improved with rapid CT diagnosis, early aneurysm treatment, and dedicated neuro-intensive care and stroke units. Survivors of rupture may make substantial gains through rehabilitation, but fatigue, memory problems, and emotional difficulties are common and may persist for years. Long-term follow-up, including repeat imaging, is recommended for those treated for aneurysms.

\section*{Prevention and Self-Care}
\begin{itemize}
    \item Control blood pressure through a healthy diet, regular physical activity, weight management, and medication as prescribed
    \item Do not smoke, and avoid second-hand smoke
    \item Limit alcohol intake and avoid stimulant drugs such as cocaine
    \item Discuss screening with your doctor if two or more close relatives have had a brain aneurysm, or if you have a predisposing condition such as polycystic kidney disease
    \item If an unruptured aneurysm is detected, attend all scheduled follow-up imaging and follow advice about activity and symptom monitoring
    \item Seek urgent medical help for any sudden severe headache rather than waiting to see whether it resolves
    \item During recovery from a bleed, follow the rehabilitation programme, attend follow-up appointments, and seek help for mood or cognitive difficulties
\end{itemize}

\section*{Sources and Further Reading}
The following organisations provide reliable, evidence-based information for patients, students, and clinicians.

\begin{itemize}
    \item NHS. \textit{Brain aneurysm: symptoms, causes, and treatment information.} https://www.nhs.uk/conditions/
    \item Mayo Clinic. \textit{Brain aneurysm: overview and patient education.} https://www.mayoclinic.org/diseases-conditions/
    \item MSD Manual Professional Version. \textit{Cerebral aneurysm and subarachnoid haemorrhage: clinical guidance.} https://www.msdmanuals.com/
    \item MedlinePlus. \textit{Cerebral aneurysm: consumer health summary.} https://medlineplus.gov/
    \item National Institute for Health and Care Excellence. \textit{Clinical guidance on the management of aneurysmal subarachnoid haemorrhage.} https://www.nice.org.uk/
    \item World Health Organization. \textit{Cardiovascular and neurological disease prevention resources.} https://www.who.int/
\end{itemize}
